\documentclass[11pt,a4paper]{article}
\usepackage[margin=1.05in]{geometry}
\usepackage{amsmath,amssymb,amsthm,mathtools}
\usepackage{booktabs}
\usepackage{microtype}
\usepackage[hidelinks]{hyperref}

\theoremstyle{plain}
\newtheorem{theorem}{Theorem}[section]
\newtheorem{proposition}[theorem]{Proposition}
\newtheorem{lemma}[theorem]{Lemma}
\newtheorem{corollary}[theorem]{Corollary}
\theoremstyle{definition}
\newtheorem{definition}[theorem]{Definition}
\newtheorem{problem}{Problem}
\theoremstyle{remark}
\newtheorem{remark}[theorem]{Remark}

\DeclareMathOperator{\rad}{rad}
\newcommand{\rhotwo}{\rho}
\newcommand{\Zpos}{\mathbb{Z}_{>0}}

\title{Least Realizers and Occupation in an Aligned\\ Two-Block Collatz Family}
\author{Elias De Jes\'us\\ \small ORCID 0009-0007-0190-9143}
\date{19 November 2026}

\begin{document}
\maketitle

\begin{abstract}
Let $T$ be the Collatz shortcut map on $\Zpos$. Among the $64$ residue classes
modulo $64$ exactly two carry the aligned six-step parity blocks
$A=111100$ and $B=111010$, namely $15$ and $7$, and on those classes six steps
act as the affine maps $F_A(x)=(81x+65)/64$ and $F_B(x)=(81x+73)/64$. We ask
whether one fixed positive integer can follow these two blocks forever.
Every finite word in $\{A,B\}$ has infinitely many positive realizers --- they
form one class modulo $64^q$ --- so the question is not about the existence of
realizers but about their size: writing $r_{\min}(q)$ for the least positive
integer completing $q$ blocks, a fixed seed exits if and only if
$r_{\min}(q)\to\infty$.

Our main tool is an \emph{occupation theorem}: if $b$ positive integers below
$H$ complete $t$ blocks, then
$r_{\min}(t+b)\ge (H+\tfrac{73}{17})\rho^{-b}-\tfrac{73}{17}$ with
$\rho=81/64$, because a seed completing $t+b$ blocks contributes $b+1$ distinct
$t$-block starts at its first $b+1$ boundaries. With the cutoff
$H_t=2\cdot 32^t$ and $B_t$ the number of $t$-block starts below $H_t$, it
follows that
$\limsup_{t\to\infty}\bigl(t\log_\rho 32-B_t\bigr)=+\infty$ already implies
$r_{\min}(q)\to\infty$; an arbitrarily slowly growing, unbounded deficit along
one subsequence suffices.

Holding the origin fixed and averaging over the \emph{depth} instead gives a
second theorem: for every fixed $\ell$ one has $D_m(0)\le 2^{5m-\ell}$ outside a
set of upper density at most $3\cdot2^{-2\ell}$, so for every $\varepsilon>0$ the
set $\{m:D_m(0)>\varepsilon\,2^{5m}\}$ has density zero --- a vanishing
\emph{fraction} of the $2^{5m}$ candidate prefixes, which still permits an
absolute count far above $84m$.

As the computational result we determine the exact least realizer at depth
$49$: it is a $75$-digit integer, $r_{\min}(49)=1.65388\ldots\times 10^{74}$,
recorded in full in Theorem~\ref{thm:min49}, whose terminal state is
$31\bmod 64$, so it completes exactly $49$ blocks and then exits. The search is
exhaustive over all $2^{49}$ words and was confirmed by a second, independent
computation. Finally we isolate the one
arithmetic estimate still missing. Writing $D_m(0)=N_{6m}(64^{5m})$, a survivor
forces $D_m(0)\ge (5\log_\rho 64)\,m-\log_\rho(x_0+\tfrac{73}{17})$ with
$5\log_\rho 64=88.274237898\ldots$, so a one-sided bound $D_m(0)\le 84m$ along
infinitely many $m$ would settle the question. The average of $D_m(c)$ over all
$64^m$ translates $c$ is exactly $1$, but the bin at the origin --- the only one
that matters --- is not controlled by that average.
\end{abstract}

\section{The family, and why long words do not settle the question}

\subsection{The shortcut map and the two aligned blocks}

Let $T\colon\Zpos\to\Zpos$ be the Collatz shortcut map,
\[
T(x)=\begin{cases} (3x+1)/2, & x \text{ odd},\\ x/2, & x \text{ even}.\end{cases}
\]
Iterating $T$ from $x$ we record the parity of the state \emph{before} each
step; six steps produce a word in $\{0,1\}^6$ and a terminal state.

\begin{lemma}\label{lem:blocks}
Among the residue classes modulo $64$, the parity word $A=111100$ occurs
exactly on $15 \bmod 64$ and the parity word $B=111010$ occurs exactly on
$7\bmod 64$. On these classes six steps of $T$ act as
\[
F_A(x)=\frac{81x+65}{64}\quad (x\equiv 15),\qquad
F_B(x)=\frac{81x+73}{64}\quad (x\equiv 7),
\]
and in both cases the numerator is divisible by $64$. We write
$c_A=65$, $c_B=73$ and $\rho=81/64$.
\end{lemma}

\begin{proof}
Each block has four odd steps and two even steps, so six steps multiply by
$3^4=81$ and divide by $2^6=64$; the additive constant is determined by the
positions of the odd steps. The residue claim is a finite check over the $64$
classes, and the maps were verified by direct six-step simulation on the first
$2000$ members of each class (\texttt{note\_checks.py}, \S1). For instance
$15\mapsto 23\mapsto 35\mapsto 53\mapsto 80\mapsto 40\mapsto 20$ with parities
$111100$ and $(81\cdot 15+65)/64=20$, and
$7\mapsto 11\mapsto 17\mapsto 26\mapsto 13\mapsto 20\mapsto 10$ with parities
$111010$ and $(81\cdot 7+73)/64=10$.
\end{proof}

Note that $20$ and $10$ are even. Completing a block says nothing about
starting the next one: the terminal parity is unrestricted, and a seed
continues only if its new state again lies in $15$ or $7 \bmod 64$. We call
such states \emph{admissible boundaries}; they are odd, and on any infinite
path every boundary is odd.

\begin{definition}
A positive integer $x$ \emph{realizes} the word $w=w_0w_1\cdots w_{q-1}\in\{A,B\}^q$
if the boundaries $x_0=x$, $x_{j+1}=F_{w_j}(x_j)$ are all defined, i.e.\ if
$x_j$ lies in the class prescribed by $w_j$ for every $j<q$. It
\emph{completes $q$ blocks} if it realizes some word of length $q$.
\end{definition}

\subsection{One class per word}

Iterating $64x_{j+1}=81x_j+c_{w_j}$ gives the telescoped identity
\begin{equation}\label{eq:carry}
64^{q}x_q=81^{q}x_0+C_w,\qquad
C_w=\sum_{i=0}^{q-1}81^{\,q-1-i}64^{\,i}c_{w_i}.
\end{equation}

\begin{lemma}\label{lem:class}
For each $w\in\{A,B\}^q$ the set of positive integers realizing $w$ is exactly
one arithmetic progression,
\[
\{x>0:\ x\equiv r_w \pmod{64^{q}}\},\qquad
r_w\equiv -81^{-q}C_w \pmod{64^{q}},
\]
where $r_w$ denotes the representative in $[0,64^q)$. Moreover
$r_w\equiv 15$ or $7 \bmod 64$, so $r_w>0$ and $r_w$ is itself the least
positive realizer of $w$.
\end{lemma}

\begin{proof}
Induction on $q$. For $q=1$ the statement is Lemma~\ref{lem:blocks}. Assume it
for $q-1$. If $x$ realizes $w$ then $x\equiv r_{w_0}\bmod 64$, and
$x_1=(81x+c_{w_0})/64$ realizes $w_1\cdots w_{q-1}$, hence lies in a single
class modulo $64^{q-1}$. On the class of $r_{w_0}$ modulo $64^{q}$ the map
$x\mapsto (81x+c_{w_0})/64$ is a bijection onto a class modulo $64^{q-1}$,
because $81$ is a unit modulo $64^{q-1}$; pulling back the class of
$x_1$ gives one class modulo $64^{q}$. Conversely every member of that class
realizes $w$, since the labels read off at each step depend only on $x$ modulo
$64^{q}$. Reducing \eqref{eq:carry} modulo $64^q$ gives the formula for $r_w$.
Finally $r_w\equiv r_{w_0}\bmod 64$ is $15$ or $7$, hence nonzero.
\end{proof}

For $q\le 3$ we also verified Lemma~\ref{lem:class} by brute force over all
$64^q$ residues and all $2^q$ words (\texttt{note\_checks.py}, \S1).

\begin{definition}
$r_{\min}(q)=\min\{r_w:\ w\in\{A,B\}^q\}$ is the least positive integer
completing $q$ blocks, and for $H>0$
\[
N_t(H)=\#\{x\in\Zpos:\ x<H,\ x \text{ completes at least } t \text{ blocks}\}.
\]
\end{definition}

Since a seed completing $q+1$ blocks completes $q$ blocks, $r_{\min}$ is
nondecreasing and $N_t(H)$ is nonincreasing in $t$.

\subsection{Why arbitrarily long words do not answer the question}

The question is open in general, not only for integers: L\'opez and Stoll
\cite{LopezStoll2009}, who computed $\Phi$ on mechanical words of irrational
slope, record in their abstract that it is unknown whether any aperiodic parity
vector $v$ has an eventually periodic --- equivalently rational --- value
$\Phi(v)$. The two-block family restricts the question to one explicit
language and to positive integer seeds. Any positive integer surviving forever
would have an aperiodic parity vector, since its block boundaries strictly
increase and an eventually periodic parity tail would force a repeated state.

By Lemma~\ref{lem:class} \emph{every} word of every length has infinitely many
positive realizers. The two-block family therefore never runs out of finite
solutions, and no counting statement about words can decide the fixed-seed
question. Equivalently: the nested classes of a fixed infinite word
$w_0w_1\cdots$ intersect in exactly one $2$-adic integer, so the symbolic model
always has survivors --- a Cantor set of them --- and the only obstruction to a
survivor in $\Zpos$ is archimedean. A fixed integer $x$ survives iff it lies in
all of the shrinking classes, and the moduli $64^q$ grow, so the question is
whether the least positive element of those classes must grow too.

\begin{proposition}\label{prop:equiv}
The following are equivalent.
\begin{enumerate}
\item No fixed positive integer completes $q$ blocks for every $q$
      (\emph{fixed-seed exit}).
\item $r_{\min}(q)\to\infty$ as $q\to\infty$.
\end{enumerate}
\end{proposition}

\begin{proof}
If $r_{\min}(q)\to\infty$ and $x$ completed every $q$, then $x\ge r_{\min}(q)$
for all $q$, which is absurd. Conversely, suppose $r_{\min}$ does not tend to
infinity. It is nondecreasing, hence bounded by some $R$. For each $q$ choose a
minimizer $x_q\le R$; as there are finitely many integers below $R$, one value
$x^\ast$ occurs for infinitely many $q$. Completing $q$ blocks implies
completing every smaller number of blocks, so $x^\ast$ completes every $q$ and
is a survivor.
\end{proof}

\subsection{Relation to the literature}

The parity-vector congruence of Lemma~\ref{lem:class} is the two-block
specialization of Terras's classical observation that a parity vector of length
$n$ determines one class modulo $2^n$ \cite{Terras1976}; see
\cite{Lagarias1985} for the general background.

L\'opez and Stoll \cite{LopezStoll2009} studied the $3x+1$ conjugacy map $\Phi$
on Sturmian parity words. Every Sturmian parity word has an irrational $2$-adic
seed; \cite[\S1.2]{DeJesusSturmian2026} traces the independent proofs of that
statement. The same paper gives a quantitative approximation bound and proves
transcendence for the characteristic Sturmian words, the two mechanical words
of intercept $0$ and all of their shifts, for every irrational slope below
$\gamma^{\ast}=(3+\sqrt5)/(4\log_2 3)=0.8258977\ldots$ --- a range that
contains the critical slope $\log_3 2$ (Theorem 1.2 and Corollary 1.4 there).
In the present two-block family each aligned six-bit block carries four ones,
so every infinite parity word of the family has rational frequency $2/3$ and
lies outside the Sturmian class. Those results are background here; the proofs
below depend on none of them.

A hypothetical survivor here has strictly increasing boundaries
($F_D(x)>x$ for admissible $x>0$), hence an infinite orbit, and its parity
vector has exactly four ones in each aligned six-step block --- an unaligned
window need not, as $001110$ occurring in $AB$ at offset $4$ shows --- so the
limiting density of ones is
\[
\tfrac{2}{3}=0.6\overline{6}>\frac{\ln 2}{\ln 3}=0.630929\ldots
\]
Monks and Yazinski \cite[Theorem 2.7(b)]{MonksYazinski2004} proved the lower
half of the corresponding density statement: for $\zeta\in\mathbb{Q}_{\mathrm{odd}}$
with a divergent orbit, $\liminf(h/\ell)\ge\ln2/\ln3$. That inequality does not
reach our family, whose density $2/3$ lies \emph{above} the critical value.
L\'opez and Stoll \cite[Theorem 1]{LopezStoll2021} state the two-sided form:
\emph{if the trajectory of some $\zeta\in\mathbb{Q}_{\mathrm{odd}}$ is divergent
(that is, if the orbit is an infinite set), then}
\[
\liminf_{\ell\to\infty}\frac{h}{\ell}=\frac{\ln 2}{\ln 3},
\]
where $h$ counts ones among the first $\ell$ digits of the parity vector,
$\mathbb{Q}_{\mathrm{odd}}$ is the ring of rationals with odd denominator, and
the lower limit is written $\lim$ in the preprint's own notation.
A positive integer lies in $\mathbb{Q}_{\mathrm{odd}}$, so fixed-seed exit for
our family is a consequence of that theorem; note that the weaker constraint
stated before it there, which excludes densities \emph{below} $\ln2/\ln3$, does
not apply, since $2/3$ lies above. Reference \cite{LopezStoll2021} is an
unrefereed preprint, and in this author's reading its \S6 argument passes from
the convergence of a real series $\Phi_{\mathbb{R}}$ to a conclusion about the
$2$-adic conjugacy $\Phi$ at a point where the two objects differ; we do not
rely on it here. The purpose of the present note is an independent and
elementary route to this consequence for aligned positive integers, and the
identification of the single estimate that such a route still needs.

\section{The occupation theorem}

\begin{lemma}[growth]\label{lem:growth}
Put $Y(x)=17x+73$. For admissible $x$,
\[
64\,Y(F_B(x))=81\,Y(x),\qquad 64\,Y(F_A(x))=81\,Y(x)-136 .
\]
Consequently $Y(x_{j+1})\le\rho\,Y(x_j)$ along any realized path, with equality
exactly at $B$ blocks, and
\begin{equation}\label{eq:growth}
x_j+\tfrac{73}{17}\ \le\ \rho^{\,j}\bigl(x_0+\tfrac{73}{17}\bigr).
\end{equation}
\end{lemma}

\begin{proof}
From $64x'=81x+c$ we get $64(17x'+73)=81(17x+73)+17c-17\cdot 73$, which is
$81Y(x)$ for $c=73$ and $81Y(x)-136$ for $c=65$. Dividing \eqref{eq:growth} by
$17$ and iterating gives the display. The identities were checked exactly on
admissible states including the two seeds of \S3 (\texttt{note\_checks.py},
\S2).
\end{proof}

\begin{theorem}[occupation]\label{thm:occ}
Let $t\ge 1$, $H>0$ and $b=N_t(H)$. Then
\begin{equation}\label{eq:occ}
r_{\min}(t+b)\ \ge\ \Bigl(H+\tfrac{73}{17}\Bigr)\rho^{-b}-\tfrac{73}{17}.
\end{equation}
\end{theorem}

\begin{proof}
If no positive integer completes $t+b$ blocks, then $r_{\min}(t+b)=\infty$ and
there is nothing to prove. Otherwise let $x_0$ complete $t+b$ blocks, with
boundaries $x_0<x_1<\cdots<x_b$ (strictly increasing, since $F_D(x)>x$ for
admissible $x>0$). For $0\le j\le b$ the state $x_j$ completes $t+b-j\ge t$
blocks, so $x_0,\dots,x_b$ are $b+1$ distinct integers each completing at least
$t$ blocks. If $x_b<H$ then all $b+1$ of them are $<H$, contradicting
$N_t(H)=b$, which counts the integers \emph{strictly} below $H$. Hence
$x_b\ge H$, and \eqref{eq:growth} with $j=b$ gives
$H+\tfrac{73}{17}\le x_b+\tfrac{73}{17}\le\rho^{b}(x_0+\tfrac{73}{17})$.
Rearranging and minimizing over $x_0$ yields \eqref{eq:occ}.
\end{proof}

Three features of the statement deserve emphasis. If $b=0$ then
\eqref{eq:occ} reads $r_{\min}(t)\ge H$, which is the definition of $N_t(H)=0$,
so the theorem interpolates between that trivial case and large $b$. The
strictness of the cutoff in $N_t(H)$ is used exactly once, to exclude
$x_b<H$. And the conclusion concerns depth $t+b$, not depth $t$: occupation
below $H$ at depth $t$ buys a floor only after $b$ further blocks, which is
precisely the trade the next corollary exploits.

\begin{definition}
$H_t=2\cdot 32^t=2^{5t+1}$ and $B_t=N_t(H_t)$.
\end{definition}

\begin{corollary}\label{cor:suff}
If
\[
\limsup_{t\to\infty}\bigl(t\log_\rho 32-B_t\bigr)=+\infty,
\]
then $r_{\min}(q)\to\infty$, i.e.\ fixed-seed exit holds.
\end{corollary}

\begin{proof}
Theorem~\ref{thm:occ} with $H=H_t$ and $32^t=\rho^{\,t\log_\rho 32}$ gives
\[
r_{\min}(t+B_t)\ \ge\ 2\rho^{\,t\log_\rho 32-B_t}-\tfrac{73}{17}.
\]
Along a subsequence on which the exponent tends to $+\infty$ the right-hand
side is unbounded, and $t+B_t\to\infty$; since $r_{\min}$ is nondecreasing,
$r_{\min}(q)\to\infty$. Proposition~\ref{prop:equiv} converts this into
fixed-seed exit.
\end{proof}

\begin{remark}
No rate is required. Because $r_{\min}$ is monotone, an unbounded deficit along
one sparse subsequence propagates to the full limit, and a deficit growing as
slowly as $\log\log t$ is enough --- it yields a floor
$2\rho^{\log\log t}$, which is unbounded although it grows more slowly than any
power of $t$. Two convenient strengthenings: a linear deficit, i.e.
$B_t\le 14t$ for arbitrarily large $t$ --- admissible since
$\log_\rho 32=14.712372983\ldots$, certified by $3^{56}<2^{89}$ --- gives the
exponential floor $2\,(2^{89}/3^{56})^{t}-\tfrac{73}{17}
=2\cdot 1.1827129\ldots^{\,t}-\tfrac{73}{17}$, the additive constant being the
one carried by Theorem~\ref{thm:occ}; and any
finite value of $\liminf_T T^{-1}\sum_{t\le T}B_t$ gives infinitely many $t$
with $B_t=O(1)$, hence a floor of order $32^t$ at depth $t+O(1)$.
\end{remark}

\section{The sharp finite result: depth 49}

\begin{theorem}\label{thm:min49}
The least positive integer completing $49$ blocks is the $75$-digit integer
\[
r_{\min}(49)=
\begin{aligned}[t]
&1653883058860887846303845081\\
&\quad 3920236748721560039331093621\\
&\quad\quad 4638366216510132943 .
\end{aligned}
\]
Its word is
\[
\texttt{AABAAABABAAAAAABABABAABBAABBABBBBBABABBBABABABBAA},
\]
and its terminal state is
\begin{align*}
&1703930977523477571221341143505225157486881936\\
&\qquad 2776884297752061966903500567997727\ \equiv\ 31 \pmod{64},
\end{align*}
which is neither $15$ nor $7$; hence $r_{\min}(49)$ completes exactly $49$
blocks and then exits.
\end{theorem}

\begin{proof}[Computation]
By Lemma~\ref{lem:class}, $r_{\min}(49)=\min_w r_w$ over the $2^{49}$ words,
where $r_w$ is the representative in $[0,64^{49})$ of
$-81^{-49}C_w$. Writing $u_i=64^{i}81^{-i-1}\bmod 64^{49}$, we have
$r_w\equiv \mathrm{base}+\sum_{i\in B(w)}s_i$ with
$\mathrm{base}=-65\sum_i u_i$ and $s_i=-8u_i$, where $B(w)$ is the set of
positions carrying $B$; all arithmetic is modulo $M=64^{49}=2^{294}$.

Split the $49$ positions into two blocks and let $\mathcal{A}$, $\mathcal{B}$
be the sets of subset sums of $s_i$ over each block, so that every word
corresponds to exactly one pair $(a,b)\in\mathcal{A}\times\mathcal{B}$ and
$r_w=(\mathrm{base}+a+b)\bmod M$. Store $\mathcal{B}$ sorted. For fixed $a$
put $v=(\mathrm{base}+a)\bmod M$; then $(v+b)\bmod M$ equals $v+b$ if
$v+b<M$ and $v+b-M$ otherwise, and since $v+b<2M$ always, the minimum over
$b\in\mathcal{B}$ is attained either at the least element of $\mathcal{B}$
(unwrapped branch) or at the least element $\ge M-v$ (wrapped branch). Taking
the better of these two candidates for every $a$ therefore examines, for each
$a$, the best $b$; as every word is some $(a,b)$, the search is exhaustive over
all $2^{49}$ words. Each candidate is found by one binary search, so the cost
is $|\mathcal{A}|$ searches in a table of size $|\mathcal{B}|$.

This was run with $|\mathcal{B}|=2^{24}$ (positions $25,\dots,48$) and
$|\mathcal{A}|=2^{25}$ (positions $0,\dots,24$): \texttt{min49.py}, $95.7$
seconds, about $1.3$\,GiB, log \texttt{min49.log}. The value above was the
running minimum after the first eighth of the search and was not improved.
\end{proof}

\paragraph{Witness versus minimality.} The two halves of
Theorem~\ref{thm:min49} are of different character and were checked
separately. That the stated integer \emph{is} a realizer is a witness
property, confirmed by replaying it block by block and, independently, by
replaying it through the $294$ individual shortcut steps with parities recorded
before each step; both routes return the same word and the same terminal state
$31\bmod 64$ (\texttt{note\_checks.py}, \S3). That it is the \emph{least} such
integer is an exhaustive claim. We re-verified it by a second computation of a
different shape: instead of searching for a minimum, \texttt{verify49.py}
counts every $49$-block start strictly below $r_{\min}(49)+1$ by an exact
cyclic-interval test, using the complementary split (stored positions
$0,\dots,23$, iterated positions $24,\dots,48$). It returns exactly one start,
the value above ($93.5$ seconds, log \texttt{verify49.log}).

\begin{proposition}[control at depth 48]\label{prop:48}
Exactly one positive integer below $2^{241}$ completes $48$ blocks, namely the
$72$-digit
\[
\begin{aligned}[t]
&1007663375116395146880531925190\\
&\quad 38392361740472623921585501397359010728775,
\end{aligned}
\]
with word \texttt{BBBBABABAAAABABBBBBBAABABABBABABBBBAAAAAABABBABB} and
terminal state
\begin{align*}
&8202715119948079618738700704624290008\\
&\qquad 326907991008520856030001878734909577507\ \equiv\ 35\pmod{64}.
\end{align*}
Consequently every positive integer below
$2^{241}=3.53\ldots\times 10^{72}$ exits before completing block $49$.
\end{proposition}

All arithmetic here is modulo $64^{48}=2^{288}$, as in the depth-49 search it
is modulo $64^{49}=2^{294}$; both implementations use that exact modulus, with
no padded precision. This is a count rather than a minimum, and it was obtained
by meet-in-the-middle over the $2^{48}$ residues in two independent
implementations --- one using the
fact that $v_2(s_i)=3+6i$ makes the late switches multiples of $64^{24}$, one
branching on the first six letters and testing an exact cyclic interval on full
residues modulo $2^{288}$ --- which agree (\texttt{note\_verifier.py},
\texttt{independent48.py}; logs in \texttt{verification/logs}). It is
consistent with and weaker than Theorem~\ref{thm:min49}, since
$r_{\min}(49)=1.65\ldots\times 10^{74}>2^{241}$; we record it because it is the
cheapest certificate of the same shape and because its single start is the
$m=8$ entry of Table~\ref{tab:bins}.

\section{The missing estimate}

\subsection{A balanced split}

Let $t=6m$ and take the cutoff $H=64^{5m}=2^{30m}=32^{t}$, which is $H_t/2$.
Set
\[
D_m(0)=N_{6m}\bigl(64^{5m}\bigr).
\]

\begin{proposition}\label{prop:nec}
If $x_0\in\Zpos$ completes every block, then for every $m\ge 1$
\begin{equation}\label{eq:nec}
D_m(0)\ \ge\ (5\log_\rho 64)\,m-\log_\rho\Bigl(x_0+\tfrac{73}{17}\Bigr),
\qquad 5\log_\rho 64=88.274237898\ldots
\end{equation}
\end{proposition}

\begin{proof}
Put $b=D_m(0)$. The survivor completes $6m+b$ blocks, so by the argument of
Theorem~\ref{thm:occ} its $b$-th boundary satisfies $x_b\ge 64^{5m}$, and
\eqref{eq:growth} gives $64^{5m}\le\rho^{b}(x_0+\tfrac{73}{17})$. Taking
$\log_\rho$ and using $\log_\rho 64^{5m}=5m\log_\rho 64$ gives \eqref{eq:nec}.
\end{proof}

\begin{corollary}\label{cor:84}
If $D_m(0)\le 84m$ for infinitely many $m$, then fixed-seed exit holds, i.e.
$r_{\min}(q)\to\infty$.
\end{corollary}

\begin{proof}
For such $m$, \eqref{eq:nec} forces
$\log_\rho(x_0+\tfrac{73}{17})\ge (5\log_\rho 64-84)m=4.2742\ldots m$, which is
unbounded, contradicting the fixedness of $x_0$. The admissibility of the
constant $84$ is the exact inequality $\rho^{84}<64^{5}$, equivalently
$3^{336}<2^{534}$, verified as an integer comparison
(\texttt{note\_checks.py}, \S2). Proposition~\ref{prop:equiv} converts the
absence of a survivor into $r_{\min}(q)\to\infty$.
\end{proof}

The constant $84$ is a benchmark only: any constant strictly below
$5\log_\rho 64=88.2742\ldots$ serves, and by the remark after
Corollary~\ref{cor:suff} so does any sublinear bound, or indeed any unbounded
deficit.

Why this cutoff is the natural one: $6m=5m+m$ splits the depth so that the
cutoff $64^{5m}$ is exactly the modulus of the first $5m$ blocks.

\begin{lemma}\label{lem:onecand}
The $2^{5m}$ words of length $5m$ have pairwise distinct classes modulo
$64^{5m}$. Hence each word $u$ has exactly one realizer $r_u$ below $64^{5m}$,
and
\[
D_m(0)=\#\bigl\{u\in\{A,B\}^{5m}:\ v_u \bmod 64^{m}\in S_m\bigr\},
\]
where $v_u$ is the terminal state of $r_u$ after $5m$ blocks and
$S_m=\{r_w\bmod 64^m: w\in\{A,B\}^m\}$, of size $2^m$.
\end{lemma}

\begin{proof}
Distinctness: if $u\ne u'$ first differ at position $k$, then
$r_u-r_{u'}\equiv\pm s_k+\sum_{i>k}\pm s_i$, and $v_2(s_i)=3+6i$, so the
difference has $2$-adic valuation exactly $3+6k<6\cdot 5m$ and is nonzero. By
Lemma~\ref{lem:class} the realizers of $u$ form one class modulo $64^{5m}$,
whose unique representative below $64^{5m}$ is $r_u>0$. Such a representative
completes $6m$ blocks precisely if its state after $5m$ blocks admits $m$
further blocks, i.e.\ if $v_u$ lies in one of the $2^m$ classes modulo $64^m$
of Lemma~\ref{lem:class} at depth $m$.
\end{proof}

\subsection{The mean over translates is one, and misses the point}

For $c\in\mathbb{Z}/64^m$ let
$D_m(c)=\#\{u\in\{A,B\}^{5m}: v_u-c \bmod 64^m\in S_m\}$, so that $D_m(0)$ is
the quantity in Lemma~\ref{lem:onecand}.

\begin{lemma}\label{lem:mean}
$\displaystyle\sum_{c\in\mathbb{Z}/64^m}D_m(c)=2^{5m}\cdot 2^{m}=64^{m}$, so
the mean of $D_m(c)$ over all $64^m$ translates is exactly $1$.
\end{lemma}

\begin{proof}
Double counting: each of the $2^{5m}$ words $u$ contributes to exactly
$|S_m|=2^m$ translates $c$, namely $c\equiv v_u-s$ for $s\in S_m$. Hence the
sum is $2^{5m}2^m=2^{6m}=64^m$, and dividing by the number of translates gives
mean $1$.
\end{proof}

This is the exact sense in which $D_m(0)\le 84m$ is "expected": $84m$ exceeds
the mean by a factor $84m$. But Lemma~\ref{lem:mean} is invariant under every
translation of the target set and therefore cannot see \emph{which} translate
is the survivor set; the bin at the origin is the one translate defined by the
same arithmetic that produces the endpoints $v_u$, and nothing in the average
constrains it. The spread is real and not small: computing the full translate
table (\texttt{note\_checks.py}, \S4) gives
\[
\max_c D_m(c)=3,\ 6,\ 8 \quad\text{for } m=1,2,3,
\]
against the mean $1$, while $D_m(0)=2,2,3$.

\begin{remark}[where the obstruction sits]
From \eqref{eq:carry} with $q=5m$,
\[
v_u \bmod 64^m=\Bigl\lfloor\frac{81^{5m}r_u+C_u}{2^{30m}}\Bigr\rfloor \bmod 2^{6m},
\qquad
\frac{65\,(81^{5m}-64^{5m})}{17}\le C_u\le\frac{73\,(81^{5m}-64^{5m})}{17},
\]
so the carry is pinned by the word and the deciding quantity is the window of
bits $30m,\dots,36m$ of the product $81^{5m}r_u$. The low $30m$ bits of that
product encode the prefix $u$ (the $2$-adic side) and the top bits are what the
cutoff controls (the archimedean side); a match is decided by the bits in
between, which neither side pins. Any proof of a bound on $D_m(0)$ must say
something about those middle bits.
\end{remark}

\subsection{Holding the origin fixed and varying the depth}

Lemma~\ref{lem:mean} averages over the $64^m$ translates at a \emph{fixed}
depth, and is blind to which translate is the survivor set. The estimate of
this subsection does the opposite: it keeps the bin at the origin and averages
over the \emph{depth}. The two are independent, and only the second says
anything about $D_m(0)$ itself.

\begin{theorem}[depth averaging]\label{thm:depthavg}
Fix an integer $\ell\ge1$ and put $s=\ell$. Then
\[
D_m(0)\ \le\ 2^{\,5m-\ell}
\]
for every $m\ge\ell$ outside a set of upper density at most $3\cdot2^{-2\ell}$.
\end{theorem}

\begin{proof}
Write $p=5m$, $n=p-\ell$, $L=64^{\ell}$, $M=64^{\ell+s}$ and $a=81^{\,n}\bmod M$.

\emph{Prefix lift.} By Lemma~\ref{lem:class} the starts below $64^{p}$ in the
class of a length-$n$ prefix $u$ are exactly $r_u+64^{n}k$ with $0\le k<L$, and
by the computation of Lemma~\ref{lem:onecand} the state after $n$ blocks is
$v_u+81^{n}k$. Since $s\le m$, a start counted by $D_m(0)$ completes in
particular $n+\ell+s$ blocks, which requires
\[
v_u+81^{n}k\in R_{\ell+s}\pmod M,\qquad\text{i.e.}\qquad
k\in a^{-1}\bigl(R_{\ell+s}-v_u\bigr)\pmod M .
\]
The admissible $k$ lie in a window of $L$ consecutive residues, and
$a^{-1}(R_{\ell+s}-v_u)$ is a rotation of $a^{-1}R_{\ell+s}$, so their number is
at most
\[
K_{\ell,s}(a):=\max_{\text{windows of }L\text{ consecutive residues}}
\#\bigl(a^{-1}R_{\ell+s}\cap\text{window}\bigr).
\]
Summing over the $2^{\,n}=2^{\,5m-\ell}$ prefixes,
$D_m(0)\le 2^{\,5m-\ell}K_{\ell,s}(81^{\,5m-\ell})$.

\emph{Phase cycle.} Since $81-1=80$ has $v_2=4$ exactly, the order of $81$
modulo $2^{j}$ is $2^{\,j-4}$ for $j\ge4$, so $81$ generates the subgroup of
units congruent to $1$ modulo $16$. Here $M=2^{6(\ell+s)}$ and that subgroup
has order $P_\ell:=2^{6(\ell+s)-4}$. The stride is $5$, which is odd, and the
order is a power of two, so $m\mapsto 81^{\,5m-\ell}\bmod M$ runs through
\emph{every} such unit, with period $P_\ell$. Consequently
$K_{\ell,s}(81^{\,5m-\ell})$ depends only on $m\bmod P_\ell$: the set of $m$ at
which it exceeds $1$ is a union of residue classes modulo $P_\ell$, so its
density exists and equals the proportion of bad units.

\emph{Pair estimate.} $K_{\ell,s}(a)\ge2$ exactly when two points of
$a^{-1}R_{\ell+s}$ lie at circular distance less than $L$. Let $\Delta$ be the
difference of two distinct elements of $R_{\ell+s}$; as in
Proposition~\ref{prop:valbound}, $v:=v_2(\Delta)=3+6i$ where $i$ is the first
index at which the two words differ. If $v\ge6\ell$ then $u\Delta$ is a nonzero
multiple of $2^{6\ell}=L$ for every unit $u$, hence at circular distance at
least $L$ from $0$. Otherwise $v\le6\ell-3$, and as $u$ runs over the units
congruent to $1$ modulo $16$ the products $u\Delta$ run uniformly over a coset
of spacing $2^{\,v+4}$ in $\mathbb{Z}/M$. That coset meets the interval
$(-L,L)$ in at most $(2L-2)/2^{\,v+4}+1$ points, so the proportion of units for
which this one pair is bad is at most
\[
\frac{2L}{M}+\frac{2^{\,v+4}}{M}\ \le\ \frac{4L}{M},
\qquad\text{using } 2^{\,v+4}\le 2^{\,6\ell+1}=2L .
\]

\emph{Union bound.} $R_{\ell+s}$ has $2^{\ell+s}$ elements, hence fewer than
$2^{\,2(\ell+s)-1}$ unordered pairs, and the proportion of bad units is at most
\[
2^{\,2(\ell+s)-1}\cdot\frac{6L}{M}=3\cdot2^{\,2\ell-4s},
\]
where the safe constant $6L/M$ absorbs the estimate above. With $s=\ell$ this
is $3\cdot2^{-2\ell}$. At every $m$ outside that set $K_{\ell,s}=1$, and the
prefix lift gives $D_m(0)\le2^{\,5m-\ell}$.

\emph{Small moduli and finite exceptions.} The argument uses $s\le m$ (only $m$
blocks remain beyond depth $5m$) and $\ell<5m$; with $s=\ell$ both hold as soon
as $m\ge\ell$, and $j=6(\ell+s)\ge12>4$ makes the order computation valid for
every $\ell\ge1$. The finitely many $m<\ell$ are absorbed into the exceptional
set without changing its density.
\end{proof}

\begin{corollary}\label{cor:dens}
For every $\varepsilon>0$,
$\operatorname{dens}\{m:D_m(0)>\varepsilon\,2^{5m}\}=0$.
\end{corollary}

\begin{proof}
Given $\varepsilon$, choose $\ell_0$ with $2^{-\ell_0}\le\varepsilon$. For
\emph{every} $\ell\ge\ell_0$,
$\{m:D_m(0)>\varepsilon2^{5m}\}\subseteq\{m:D_m(0)>2^{\,5m-\ell}\}$, which by
Theorem~\ref{thm:depthavg} has upper density at most $3\cdot2^{-2\ell}$.
Letting $\ell\to\infty$ gives upper density $0$. The order of quantifiers is:
$\varepsilon$ first, then $\ell$ depending on $\varepsilon$ alone, then $m$.
\end{proof}

\begin{remark}
A vanishing \emph{fraction} is not a small \emph{count}: $2^{\,5m-\ell}$ exceeds
$84m$ for every fixed $\ell$ once $m$ is large, so Theorem~\ref{thm:depthavg}
bounds the proportion of surviving prefixes and leaves the absolute count
exponentially far from the target of Corollary~\ref{cor:84}.
\end{remark}

\begin{remark}
The union bound is loose. An exact census of the phases
(\texttt{averaging.py}; see Appendix~\ref{app:pkg}) gives a bad proportion of
$24/256=0.0938$ at $\ell=s=1$ against the bound $0.75$, and
$17912/1048576=0.0171$ at $\ell=s=2$ against $0.1875$ --- a factor $8$ to $11$.
The method also survives past its nominal range: at $(\ell,s)=(3,1)$ the union
bound is vacuous, yet $23.8\%$ of depths are still good.
\end{remark}

\subsection{What first-difference valuations alone can prove}

\begin{proposition}\label{prop:valbound}
$B_t\le 2^{K}$ with $K=\lceil (5t-2)/6\rceil$: a $t$-block start below
$H_t=2^{5t+1}$ is determined by the first $K$ letters of its word.
\end{proposition}

\begin{proof}
If two starts below $H_t$ have words first differing at position $k$, then as
in Lemma~\ref{lem:onecand} their difference is a nonzero multiple of
$2^{3+6k}$, while both lie in $[0,2^{5t+1})$, so $2^{3+6k}<2^{5t+1}$ and
$k<(5t-2)/6$. Hence distinct starts have distinct prefixes of length $K$.
\end{proof}

Refining by prefix length $k$ gives $B_t\le 2^{5t-2-5k}+2^{k}$, optimized at
$k\approx(5t-2)/6$ as well, so $2^{5t/6}$ is the ceiling of this argument and
not an artifact of one proof. The ceiling is attained by a family sharing all
the data the argument uses.

\begin{proposition}\label{prop:family}
Let $\mathcal{F}_t=\{7+8\sum_{i<t}\epsilon_i 64^{i}:\epsilon\in\{0,1\}^t\}$.
Then $\mathcal{F}_t$ has the admissible residues $7,15\bmod 64$, is nested in
the same sense (a prefix of the $\epsilon_i$ fixes the value modulo $64^k$),
has first-difference valuations $v_2=3+6i$, and contains exactly
$2^{\lceil (5t-2)/6\rceil}$ members below $2^{5t+1}$.
\end{proposition}

\begin{proof}
Since $8\cdot 64^{i}=2^{3+6i}$, the members below $2^{5t+1}$ are precisely
those with $\epsilon_i=0$ whenever $3+6i\ge 5t+1$: any single such term already
exceeds the cutoff, while the sum of all terms with $3+6i\le 5t$ is at most
$7+2^{5t}\cdot\frac{64}{63}<2^{5t+1}$. Hence the count is
$2^{\#\{i:3+6i\le 5t\}}=2^{\lfloor (5t-3)/6\rfloor+1}=2^{\lceil (5t-2)/6\rceil}$,
using $\lfloor (n-1)/6\rfloor+1=\lceil n/6\rceil$.
\end{proof}

Proposition~\ref{prop:family} was verified for $t=3,\dots,14$, rounding
included, and the floor/ceiling identity for $t=3,\dots,59$
(\texttt{note\_checks.py}, \S4). It follows that admissible residues, nesting
and first-difference valuations together cannot give any bound better than
$2^{5t/6}$ --- in particular they cannot give $84m$ or $14t$ --- so any proof of
the missing estimate must use arithmetic specific to the block maps.

\subsection{The finite counts}

Table~\ref{tab:bins} lists $D_m(0)$ for $m\le 8$; Table~\ref{tab:Bt} lists
$B_t$ for $t\le 40$. Every value is far below the respective allowance
($84m$, $14t$), and both sequences hover around their means ($1$ and $2$).
These are finite counts and are evidence only; they say nothing about
unbounded $m$ or $t$.

\begin{table}[ht]
\centering
\begin{tabular}{lrrrrrrrr}
\toprule
$m$ & 1 & 2 & 3 & 4 & 5 & 6 & 7 & 8\\
$D_m(0)$ & 2 & 2 & 3 & 0 & 0 & 1 & 2 & 1\\
allowance $84m$ & 84 & 168 & 252 & 336 & 420 & 504 & 588 & 672\\
\bottomrule
\end{tabular}
\caption{The distinguished bin. $m\le 6$ from \texttt{note\_checks.py};
$m=7$ from \texttt{m7.py}; $m=8$ from Proposition~\ref{prop:48}, whose unique
start below $2^{241}$ is in fact below $2^{240}=64^{40}$.}
\label{tab:bins}
\end{table}

\begin{table}[ht]
\centering\footnotesize
\setlength{\tabcolsep}{3.1pt}
\begin{tabular}{l*{19}{r}}
\toprule
$t$   & 3&4&5&6&7&8&9&10&11&12&13&14&15&16&17&18&19&20&21\\
$B_t$ & 2&2&1&2&2&4&1&2&1&6&0&1&2&2&1&4&2&0&0\\
\midrule
$t$   & 22&23&24&25&26&27&28&29&30&31&32&33&34&35&36&37&38&39&40\\
$B_t$ & 3&0&1&1&1&1&3&2&0&1&2&2&3&2&3&2&1&1&1\\
\bottomrule
\end{tabular}
\caption{$B_t=N_t(2\cdot 32^t)$ for $3\le t\le 40$ (\texttt{windowcount.py}).
The sum is $65$, the mean $65/38=1.7105\ldots$ against the translate mean $2$,
and the maximum is $6$ at $t=12$ against the allowance $14\cdot 12=168$. Also
$B_{48}=1$.}
\label{tab:Bt}
\end{table}

\section{A second expression for the same count}\label{sec:split}

The split used so far is $6m=5m+m$: a prefix of $5m$ blocks whose least
realizer is compared with the cutoff, and $m$ further blocks. A balanced split
$6m=3m+3m$ expresses the same number as a count of ordered pairs, and converts
it into a correlation between two explicit lists.

Throughout this section $n=3m$, $B=64^{\,3m}$, $Q=B^{2}=64^{\,6m}$,
$W=64^{\,2m}$, $H=BW=64^{\,5m}$ and $J=6m$, so that $B/W=2^{J}$. Write
$[z]_M$ for the least nonnegative residue, put $u_i=[64^{i}81^{-i-1}]_Q$ for
$0\le i<2n$, and for half-words $s,t\in\{A,B\}^{n}$ set
\[
P_s=\Bigl[-\sum_{i<n}c_{s_i}u_i\Bigr]_Q,\qquad
T_t=\Bigl[\ \sum_{i<n}c_{t_i}u_{n+i}\Bigr]_Q,\qquad
a_s=\Bigl\lfloor \frac{P_s}{B}\Bigr\rfloor,\qquad b_t=\frac{T_t}{B}.
\]

\begin{proposition}[half-word reduction]\label{prop:halfword}
$B$ divides $T_t$ for every $t$; the least realizer of the concatenated word
$st$ satisfies
\[
r_{st}=(P_s\bmod B)+B\,[\,a_s-b_t\,]_B ;
\]
and
\[
D_m(0)=\#\bigl\{(s,t)\in\{A,B\}^{n}\times\{A,B\}^{n}\ :\ [\,a_s-b_t\,]_B<W\bigr\},
\]
counted as ordered pairs, the values $a_s$ being taken with multiplicity.
\end{proposition}

\begin{proof}
$v_2(u_{n+i})=6(n+i)\ge6n$, so $B=2^{6n}$ divides each $u_{n+i}$ and hence
$T_t$; thus $b_t$ is an integer. By Lemma~\ref{lem:class},
$r_{st}=[-\sum_{i<2n}c_{(st)_i}u_i]_Q=[P_s-T_t]_Q$. Write $P_s=Ba_s+\rho$ with
$0\le\rho<B$ and $T_t=Bb_t$; then $P_s-T_t\equiv B(a_s-b_t)+\rho$, and since
$0\le\rho<B$ and $0\le B[a_s-b_t]_B<B^{2}=Q$ the displayed expression is the
least residue. Finally $r_{st}<H=BW$ if and only if $[a_s-b_t]_B<W$, because
$\rho<B$ makes the two cases $[a_s-b_t]_B\le W-1$ and $[a_s-b_t]_B\ge W$
decide the inequality; the correspondence $st\mapsto(s,t)$ is a bijection onto
ordered pairs, and distinct half-words may share a value of $a_s$, which is why
multiplicities are retained.
\end{proof}

Partition $[0,B)$ at level $j$ into $2^{j}$ half-open bins of equal width, so
that level $J$ has width $W$. Let $A$ be the multiset $(a_s)_s$ and let $Z$
contain both $b_t$ and $[b_t+W]_B$ for every $t$. A pair counted by
Proposition~\ref{prop:halfword} has $a_s\in[b_t,b_t+W)$, which meets the bin of
$b_t$ and the next bin cyclically, so with
$C_j=\sum_{I\text{ at level }j}A(I)Z(I)$ we get $D_m(0)\le C_J$. For a bin $I$
at level $j$ with children $I_{\mathrm{left}},I_{\mathrm{right}}$ put
$\Delta A(I)=A(I_{\mathrm{left}})-A(I_{\mathrm{right}})$ and likewise for $Z$.

\begin{proposition}\label{prop:bins}
With $e_j:=2^{\,j-J}\sum_{I\text{ at level }j}\Delta A(I)\Delta Z(I)$,
\[
D_m(0)\ \le\ C_J\ =\ 2+\sum_{j=0}^{J-1}e_j .
\]
\end{proposition}

\begin{proof}
For one parent, $A(I_L)Z(I_L)+A(I_R)Z(I_R)=\tfrac12\bigl(A(I)Z(I)+\Delta
A(I)\Delta Z(I)\bigr)$, by expanding $A_{L,R}=(A\pm\Delta A)/2$. Summing over
level $j$ gives $C_{j+1}=\tfrac12(C_j+G_j)$ with
$G_j=\sum_{I}\Delta A(I)\Delta Z(I)$, so
$C_J=2^{-J}C_0+\sum_{j<J}2^{\,j-J}G_j$. Here
$C_0=|A|\,|Z|=2^{\,3m}\cdot2^{\,3m+1}=2^{\,6m+1}$ and $2^{-J}=2^{-6m}$, which
gives the constant $2$.
\end{proof}

Since $J=6m$, the test $e_j\le13$ at every level would give $C_J\le2+78m\le84m$
and hence, by Corollary~\ref{cor:84}, fixed-seed exit. That test is verified
only at the six depths of Table~\ref{tab:split}.

\paragraph{A worked example: $m=2$, $j=9$.} Here $B=2^{36}$, $W=2^{24}$,
$J=12$, $|A|=64$ and $|Z|=128$. The children sit at level $10$, in bins of
width $2^{26}$: $125$ of them contain a value of $A$ or of $Z$, and $78$
contain a value of $Z$. These lie under $110$ occupied parent bins at level
$9$. Summing the suffix imbalances over those parents gives exactly
\[
\sum_{I}\Delta Z(I)=0,
\]
yet the products $\Delta A(I)\Delta Z(I)$ are nonnegative at every parent but
one and total $G_9=17$, so $e_9=17/2^{\,12-9}=17/8$ --- the largest $e_j$ at
this depth. The thirteen contributing parents carry products
$+2,+1,+1,+1,+1,+2,+2,+2,+2,+2,+2,+1$ and $-2$.

Exact balance after summing over parents therefore coexists with a positive
correlation inside them. The mechanism is visible on each side. Pair each suffix having $A$ at
position $i$ with the suffix obtained by changing that letter to $B$, keeping
the choice of $W$-shift fixed. The paired residues differ by the \emph{exact}
integer $\tau_i=8u_{n+i}/B$ modulo $B$, with $v_2(\tau_i)=3+6i$. Whenever that
difference is an odd multiple of the child-bin width, the partners lie in bins
of opposite parity, so the two contributions cancel and $\sum_I\Delta Z(I)=0$.
The additional shift by $W$ produces the same cancellation at $j=J-1$. The
pairing is an involution, not a translation: a fixed shift by $+\tau_i$ need
not preserve $Z$. At $m=1$, where $B=262144$ and $\tau_0\equiv116232$, both
$25921$ and $25921+\tau_0=142153$ lie in $Z$ while $25921+2\tau_0=258385$ does
not. So the cancellation is of signed contributions in pairs, and
$\Delta Z(I)$ stays nonzero on many individual bins ($54$ of the $110$ parents
in the example). On the prefix side the shifts are not exact: flipping letter
$i$ of $s$ moves $P_s$ by $\mp8u_i$, and $B\nmid 8u_i$ for $i<n$, so $a_s$ moves
by the exact quotient \emph{plus} a carry determined by $P_s\bmod B$ (measured
at $m=3$: $56$, $485$ and $153$ of $4000$ flips for letters $0,1,2$). The
unresolved quantity is therefore the relative placement of the prefix residues
$a_s$ and the suffix residues $b_t$, that is the correlation
$\sum_I\Delta A(I)\Delta Z(I)$ itself; the marginal identities above leave it
untouched.

\section{Two open problems}

\begin{problem}[equivalent to fixed-seed exit]
Give an independent proof that $r_{\min}(q)\to\infty$ for the aligned two-block
family; equivalently, that no fixed positive integer follows $A$ and $B$
forever.
\end{problem}

By Proposition~\ref{prop:equiv} the two formulations are the same statement.
This restricted consequence of \cite[Theorem 1]{LopezStoll2021} is what
motivates the paper: the preprint's theorem would settle it, but its proof is
one we do not rely on, and the lower-density inequality of
\cite[Theorem 2.7(b)]{MonksYazinski2004} does not reach density $2/3$. What the
paper proves independently is everything else: the occupation theorem
(Theorem~\ref{thm:occ}) and its sufficient conditions, the exact finite
certificates of \S3, the depth-averaging theorem
(Theorem~\ref{thm:depthavg}) and its density corollary, the half-word
reduction (Proposition~\ref{prop:halfword}), and the valuation ceiling
(Propositions~\ref{prop:valbound} and \ref{prop:family}).

\begin{problem}[sufficient: a one-sided bound on the distinguished bin]
Prove that $D_m(0)\le 84m$ for infinitely many $m$ --- or, more weakly, that
$5m\log_\rho 64-D_m(0)$ is unbounded above.
\end{problem}

By Corollary~\ref{cor:84} the first assertion suffices for Problem 1; by the
remark after Corollary~\ref{cor:suff} so does the second --- an unbounded
deficit below $5m\log_{81/64}64=88.274237898\ldots m$ --- which is the weakest
form we know. Theorem~\ref{thm:depthavg} does not settle either: it bounds the
proportion $D_m(0)/2^{5m}$, not the count. Problem 2 is a one-sided counting estimate at the single bin
$c=0$, with a factor $84m$ of slack over the exact mean of
Lemma~\ref{lem:mean}; what it is not is a statement about the average, which is
already known exactly and is insensitive to the bin in question.

\appendix
\section{Verification package}\label{app:pkg}

The verification scripts, reference logs, and instructions for reproducing the
computations are available at
\url{https://github.com/innerlightr-wq/collatz-two-block-occupation}.

All computations are exact integer or rational arithmetic in Python~3 with no
external dependencies, no network access, and fixed outputs. Logs accompany
each script in \texttt{verification/logs}.

Table~\ref{tab:split} records the diagnostics of \S\ref{sec:split}. The
exact counts $D_m(0)$ agree with Table~\ref{tab:bins}, which carries them
through $m=8$; the bin bound $C_J$ and the level maxima are computed from the
half-word lists. The sufficient test $e_j\le13$ holds at all six depths, with a
margin of about six, and is verified nowhere else.

\begin{table}[ht]
\centering
\begin{tabular}{lrrrrrr}
\toprule
$m$ & 1 & 2 & 3 & 4 & 5 & 6\\
\midrule
$D_m(0)$ & 2 & 2 & 3 & 0 & 0 & 1\\
$C_J=2+\sum_j e_j$ & 3 & 6 & 5 & 1 & 0 & 3\\
$\max_j e_j$ & $1$ & $17/8$ & $5/4$ & $1/2$ & $7/16$ & $1$\\
\bottomrule
\end{tabular}
\caption{The half-word diagnostics of \S\ref{sec:split}
(\texttt{seedcount.py}). $J=6m$; the bound of Proposition~\ref{prop:bins} is
$D_m(0)\le C_J$ at each depth.}
\label{tab:split}
\end{table}

\begin{center}\small
\begin{tabular}{lllr}
\toprule
script & what it certifies & checks & time\\
\midrule
\texttt{note\_checks.py} & every numerical claim in this note & 37912 & 70\,s\\
\texttt{min49.py} & $r_{\min}(49)$, Theorem~\ref{thm:min49} & --- & 96\,s\\
\texttt{verify49.py} & minimality, complementary split & --- & 94\,s\\
\texttt{windowcount.py} & $B_t$ census, Proposition~\ref{prop:valbound} & 110 & 55\,s\\
\texttt{prefix6m.py} & $D_m(0)$, Lemma~\ref{lem:onecand}, \ref{prop:family} & 34314 & 95\,s\\
\texttt{m7.py} & $D_7(0)=2$ & --- & 60\,s\\
\texttt{averaging.py} & Theorem~\ref{thm:depthavg} & 1046 & 55\,s\\
\texttt{seedcount.py} & Prop.~\ref{prop:halfword}, \ref{prop:bins}, Table~\ref{tab:split} & 4204 & 200\,s\\
\texttt{rev2\_checks.py} & the claims new in this revision & --- & 30\,s\\
\texttt{note\_verifier.py} & Proposition~\ref{prop:48} & --- & 2\,min\\
\texttt{independent48.py} & Proposition~\ref{prop:48}, independently & --- & 8\,min\\
\bottomrule
\end{tabular}
\end{center}

The depth-$48$ and depth-$49$ runs store $2^{24}$ multi-precision integers and
need about $1.3$\,GiB of memory; everything else runs in well under $1$\,GiB.

\begin{thebibliography}{9}
\raggedright
\bibitem{Lagarias1985} J.~C. Lagarias, \emph{The $3x+1$ problem and its
generalizations}, Amer. Math. Monthly \textbf{92} (1985), 3--23.
\textsc{doi:}\,10.2307/2322189.
\bibitem{DeJesusSturmian2026} E. De Jes\'us, \emph{Transcendence of the $3x+1$
conjugacy map on Sturmian words}, version~2, 5 November 2026,
\url{https://hexagonmath.org/2611.00108v2};
\textsc{doi:}\,10.5281/zenodo.24015456.
\bibitem{LopezStoll2009} J. L\'opez and P. Stoll, \emph{The $3x+1$ conjugacy
map over a Sturmian word}, Integers \textbf{9} (2009), \#A13, 141--162.
\textsc{doi:}\,10.1515/INTEG.2009.014.
\bibitem{MonksYazinski2004} K.~G. Monks and J. Yazinski, \emph{The
autoconjugacy of the $3x+1$ function}, Discrete Math. \textbf{275} (2004),
219--236. \textsc{doi:}\,10.1016/S0012-365X(03)00125-0.
\bibitem{LopezStoll2021} J. L\'opez and P. Stoll, \emph{The $3x+1$ periodicity
conjecture in $\mathbb{R}$}, preprint, arXiv:2101.12747v1 [math.NT], 29 January
2021. \textsc{doi:}\,10.48550/arXiv.2101.12747.
\bibitem{Terras1976} R. Terras, \emph{A stopping time problem on the positive
integers}, Acta Arith. \textbf{30} (1976), 241--252.
\textsc{doi:}\,10.4064/aa-30-3-241-252.

\end{thebibliography}

\end{document}
